\section{Conclusion}
\label{sec:conclusion}

An agent harness lets a model act. Someone still has to decide what is worth
attempting, whether the attempt succeeded, and when the answer belongs to a person.
That seat---the Driver---has been occupied by a human in every deployed agent system,
which is why agent progress has been rate-limited by human waking hours rather than
by model capability. \systemname{} is our attempt to hand that seat over.

The enabling condition turned out not to be a stronger model. It was giving the
completion decision an owner other than the executor. A Reviewer that runs read-only
cannot repair the evidence it judges even when that would make the round look
better; a Reviewer that can return \code{blocked} can decline to certify rather than
manufacture a success; an authority boundary that stops for credentials, money, and
irreversible actions while absorbing ordinary technical failure keeps the escalations
rare and meaningful. These are usually described as quality features. They are more
accurately described as the load-bearing wall of unattended operation: without them,
a person has to stay in the room, because nothing else can be trusted to say when the
work is done.

The delegation is measurable, so we measured it. Across \DutyCampaigns{} campaigns,
\DutyWallHours{} wall-clock hours, and \DutyEvents{} events, the runtime asked for a
human \DutyInterruptions{} times---one every \DutyInterval{} hours---and only
\DutyJudgmentPercent\% of those requests were research judgment; the rest were broken
infrastructure, authorizations the machine should not grant itself, and missing
context. With work in hand, duty cycle runs \DutyCycleLow--\DutyCycleHigh\%, and the
worst campaign we observed was throttled by a failed GPU driver rather than by
hesitation. The unattended hours produced artifacts rather than uptime, most sharply in the
domain with the least forgiving verifier. Section~\ref{sec:infra-vertical} records
the detail: a FlashAttention-4 adaptation landing within
\FlashDAGapLow--\FlashDAGapHigh\% of the native kernel for under
\FlashDACostCNY{} CNY; an SGLang integration finished after more than
\SGLangHumanHours{} hours of human effort had not finished it; and \FLAPRTotal{}
kernels accepted by outside maintainers with no blocking changes requested, one of
them a correctness fix with no speedup to report because before it the test suite
could not run at all. Across seven benchmark arenas the same runtime reaches
approximately \SWEProArgusAccuracy\% on SWE-Bench Pro against
\SWEProDirectAccuracy\% for direct Copilot, while declaring \ReviewerFirstBlocked{}
tasks blocked rather than reporting an unsupported success---the same gate seen from
the other side.

Long-horizon research also requires the Driver to do something a fixed-target executor
cannot: change the target. The objective an expert states before the work begins is a
hypothesis formed with the least information anyone will ever have, so a runtime that
treats it as authority will execute faithfully toward a destination the evidence has
already refuted. \systemname{} is evidence-driven instead: the current objective
stands while evidence supports it and yields when evidence does not, with revision
admitted only through an explicit role boundary and recorded with its justification.
That last requirement is what separates a warranted pivot from a rationalized
failure, since a rationalizing system can supply the narrative but not the refuting
measurement. Persistent campaign state, four-role authority separation, and
verification-gated updates carry that progress across model calls, tool runs,
revisions, and restarts.

The self-improvement claim separates into two axes. The parameter-invariant axis is
what this report measures and is deliberately verification-gated: candidate memories,
skills, procedures, verifiers, routing decisions, and rejected routes alter future
work only after evidence checks and an authorized commit, while model parameters stay
fixed. A failed branch thus becomes reusable progress without letting an unverified
narrative rewrite the campaign. That axis compounds, but within a ceiling---no
accumulation of state makes a frozen backbone reason past its limit. The
parameter-changing axis moves the ceiling, and reaching it requires producing
weights; \systemname{} built and ran a from-scratch 1B pretraining stack on an
eight-GPU B200 node, which is the entry condition rather than the closed loop. The
two axes are coupled: operating on the first produces the typed trajectories the
second would train on, so a runtime that self-evolves while recording how it did so
is assembling its own supervision as a byproduct of working. Training on that corpus,
and measuring whether the resulting model improves the runtime that built it, is what
remains.

The report-level $K_t/X_t$ model and \code{ManagerAdmit} operator make the
evidence-and-authority requirements for contract refinement explicit; they summarize
distributed runtime surfaces rather than one atomic API. The process-to-capability
theory explains how this runtime converts Token activity
into retained research progress. Typed trajectories contain more decision-relevant
information than final artifacts alone, review separates proposal from admission,
and reuse changes the starting point of future missions. The mathematical
campaign makes this mechanism visible through one rejected route and six accepted
theorem-frontier updates produced by the Manager, Planner, Engineer, and Reviewer.

The six-paper production case provides a second vertical view. Across
\PaperCaseCampaignHours{} aggregate campaign-hours, Argus repeatedly survives
review rejection, Stage rollback, and session replacement while producing two
AAAI- and four ACL-formatted manuscripts. In the representative campaign, seven
rejected method routes become a 4,500-row negative-results study, and two late
submission rollbacks repair the evidence package without resetting the paper. This
illustrates how the runtime can connect hypothesis revision, experiment
execution, and sustained scientific communication, while the stale assurance
snapshot in one project identifies a concrete synchronization failure for future
work.

These trajectories also define a natural interface to model training. Supervised,
preference-based, and reinforcement-learning methods can use the accumulated
objectives, actions, measurements, revisions, and state transitions to internalize
parts of the long-horizon research loop.
